\documentclass[11pt]{article}
\usepackage[T1]{fontenc}
\usepackage{lmodern}
\usepackage{amsmath,amssymb}
\usepackage{booktabs}
\usepackage{array}
\usepackage[a4paper,margin=2.2cm]{geometry}
\usepackage[hidelinks]{hyperref}
\newcommand{\doi}[1]{\href{https://doi.org/#1}{\texttt{#1}}}
\setlength{\parskip}{3pt}
\title{\vspace{-1.5cm}\textbf{Packages v4.1 --- rename, sign audit, and expanded basics}\\[2pt]
\large Milestone report}
\author{Dr.\ Vipin \\ Gurugram University}
\date{30 October 2026}
\begin{document}
\maketitle
\vspace{-8mm}
\footnotesize

\section*{1. The delivery, with counts}

Three archives, each with exactly two \LaTeX\ files. Word counts are of the source, so they run a little
above the compiled count; section counts exclude the bibliography.

\begin{center}\footnotesize
\begin{tabular}{lrrrrr}
\toprule
archive & plan words & plan sections & basics words & basics sections & symbols in table\\
\midrule
\texttt{Lamb\_Sumit.zip} & 3280 & 6 & 8840 & 23 & 41\\
\texttt{Lamb\_Neha.zip} & 2820 & 6 & 6440 & 23 & 35\\
\texttt{Lamb\_Anu.zip} & 2310 & 6 & 6040 & 23 & 34\\
\bottomrule
\end{tabular}
\end{center}

The three basics documents now stand on their own. Each opens with waves, frequency, wavelength, wavenumber,
phase velocity, dispersion, plates, modes, cutoffs, symmetric and antisymmetric motion; then complex numbers
from zero, the complex wavenumber and attenuation; then stress, strain, Hooke's law and balance laws; then the
field theory of the paper; then heat, Fourier, relaxation and MGT; then the size effect; then what a model may
assume; then a table of every symbol in the plan with its unit and a rough size; a glossary; and a closing
section on testing your own work. Neha's and Anu's documents reach the material only after the general
apparatus, as instructed.

\section*{2. The three corrections in this round}

\textbf{The name.} Every occurrence of ``modified gradient theory'' is gone from all six files. The first use
in each document now reads Moore-Gibson-Thompson (MGT), and the sentences that were built around the wrong
name were rewritten rather than patched: for example the abstract of the Sumit plan now says that the heat
equation is of the MGT type in its base form, and the Anu plan says that Matlam and Kedar applied the MGT
theory to an electro-magneto-thermoelastic medium. In the basics documents the theory has its own section,
where it is introduced by name with its two authors' dates and then called MGT for the rest of the text.

\textbf{The signs.} The reviewer was right and the file was wrong. I re-derived both coefficients from the
flux law, and the audit is in \texttt{internal/sumit\_audit.md}. Starting from
$j_i=-n b_nE_i+n b_nS_n\theta_{,i}$ and $\partial_t\Delta n=-\nabla\!\cdot\!j_i$, the carrier equation gives
$\beta_{\rm th}=-b_nn_0S_n$, positive for $n$-type, so a hot spot loses carriers and they gather where it is
cold. The heat delivered to the lattice is $-\nabla\!\cdot\!(\Pi_{\rm th}J_i)$ with $J_i=-qj_i$, which is
$-q\Pi_{\rm th}\partial_t\Delta n$; because the bracket of the heat equation carries the negative of each
source, the entry is $+q\Pi_{\rm th}\partial_t\Delta n$. Both changes are in the plan and in the model
document, and the Onsager ratio now carries its sign, $\Pi_{\rm th}/\beta_{\rm th}=-T_0/(b_nn_0)$.

The same audit was done for the other two papers and is in \texttt{internal/neha\_audit.md} and
\texttt{internal/anu\_audit.md}: every pair that shares a constant is tested for reciprocity and for
dimensions, including the phonon-phason pair $R_{ijkl}=R_{klij}$, the sign of the phason dissipation
$\gamma\dot w^2\ge0$, the magnetoelectric pair $\alpha$ appearing with the same sign in $D$ and $B$, and the
two caloric pairs. Two by-products: the unit of $\alpha$ must be s/m, which follows from either of its two
appearances, and the phason friction must be stated positive or the material becomes a pump.

\textbf{The length.} The three basics documents went from about 1500 words to 6000 to 8800, with the sections
listed in section 1 above, a symbol table of every symbol in the corresponding plan, and a closing section on
self-testing. Nothing was added that gives away a boundary condition, a solution step, or a number to compare
against.

\section*{3. What has not changed since the last review}

The seven corrections from the previous round stand: the Lim form of the nonlocal strain gradient law with the
operator on inertia and on $c{:}\varepsilon$ only, the equivalent momentum form, all piezoelectric,
piezomagnetic, thermal and phason couplings local as a stated modelling choice, the Peltier term in place of
the band-gap term, the three withdrawn novelty claims, the corrected Lubensky reference with its icosahedral
caveat, the separated symbols, and the sentence in each plan that hands the extra face conditions to the
scholar.

The claims remain narrow, and the searches behind them stand. Paper A is a free-face plate with a size effect
and attenuation curves, not a new thermoelectric constitutive law; Qu et al.\ (2023) and Gupta et al.\ (2025)
own that ground. Paper C is attenuation in a layered magneto-electro-elastic plate with MGT heat and an
effective magnetoelectric coefficient carried as a band, not the discovery of the pyroelectric or
pyromagnetic coefficient; Kondaiah et al.\ (2013) and Tocci Monaco et al.\ (2021) own those. Paper B is the
telegraph phason plus a relaxation-time heat equation, and its novelty remains provisional until the 2026
Acta Mechanica Solida Sinica paper is read.

\section*{4. What is deliberately absent, and what is left on my side}

No boundary conditions, no dispersion determinant, no solution method, no numbers to compare against. The
internal folder carries what the packages lack: three equation source logs marked [R], [S], [T], [D] or [P],
three benchmark files with every tolerance and every number we hold, three sign and unit audits, the change
log, and the style record.

No LaTeX engine exists in this workspace, so nothing was compiled here; the code is kept to the eight allowed
packages, with no images and no custom environments. If a page count comes out too long once you compile, the
cuts should come from the figure paragraphs in the plans, not from the model.

The week-one gates are unchanged. Sumit confirms equations (17) to (19) of Lim et al.\ on the page and the
Peltier term in Qu et al. Neha obtains the 2026 quasicrystal thermoelastic paper and writes down what heat
operator it uses. Anu reads the magnetoelectric sign convention in Chen, Pan and Chen and the thermal face
conditions in Kakar and Kakar. No calculations are set for the scholars yet.

\end{document}
